\thispagestyle{empty}
\vspace*{1.0cm}

\begin{center}
    \textbf{Abstract}
\end{center}

\vspace*{0.5cm}

\noindent

Monitoring plant photosynthesis across large areas is important for measuring region scale terrestrial carbon uptake and agricultural productivity. Solar-induced chlorophyll fluorescence (SIF), a faint light emission, is emitted by chlorophyll during photosynthesis and provides a remote sensing observable signal to gross primary production. NASA's OCO-2 satellite provides the highest resolution available SIF data, but it is spatially and temporally sparse. Existing state-of-the-art enhanced SIF products are limited to 500m resolution, which limits its applicability for the isolation of crop pure SIF signals in small irregularly shaped German crop fields. It is therefore necessary to construct a much higher resolution SIF product using spatially and temporally complete high resolution predictors primarily from Sentinel-2 satellite data. One unresolved issue is how model architecture and representation of predictors as either mean aggregated tabular summaries or spatial image data affect prediction accuracy and spatial consistency when evaluated within the same study framework. The influences of spatial aggregation, sample density, temporal mismatch between SIF and its predictors, and agro-climatic variation on prediction accuracy also require further evaluation. The study also examines whether crop-pure enhanced SIF provides more useful information for winter wheat yield prediction than crop-mixed SIF observations. This study modelled quality filtered combined OCO-2 SIF at 757nm and 771nm with 20m resolution Sentinel-2 spectral indices, radiation, seasonal and crop composition predictors for Germany during February-July, 2019-2024. Tree based models like Random Forest, XGBoost and neural network models like image based U-Nets were evaluated. The best U-Net model was trained on fixed-grid 4km cells containing at least five same date OCO-2 samples across nine MGRS tiles covering Germany and externally tested at native OCO-2 footprint support using 2,000 observations from two unseen tiles. The crop yield experiment compared the influence of crop-pure SIF and mixed-crop SIF across 48 Bavarian NUTS 3 regions during 2019-2022 and 2024. The Spatial Aggregate (4km) U-Net produced an internal test RMSE of 0.0701 and a $R^2$ of 0.761 while an external evaluation at the native OCO-2 footprint support for observations within the trained SIF range produced a RMSE of 0.1434 and a $R^2$ of 0.349. The external evaluation was broadly comparable with existing state-of-the-art SIF enhancement methods, although direct numerical comparison is limited by differences in datasets, geographic domains, and validation support. Native footprint CNN and tree based results showed similar performance metrics, indicating that target construction and spatial support were more influential than model architecture alone, although CNNs could better represent neighbourhood and field boundary patterns. Crop-pure SIF achieved a 2024 yield RMSE of 0.9406~$\mathrm{t\,ha^{-1}}$, compared with 1.0285 - 1.0874~$\mathrm{t\,ha^{-1}}$ for crop mixed SIF footprints, corresponding to 8.5\% and 13.5\% reduction respectively. However, the yield experiment assesses only the effect of SIF purity rather than the validity of overall yield prediction performance. The thesis provides a reproducible framework that performs OCO-2 gap filling and resolution enhancement without requiring spatially complete parent SIF support while treating its 20m outputs as model derived predictions validated at native OCO-2 footprint support. It also provides preliminary evidence that crop specific SIF isolation is useful for agricultural yield monitoring and identifies priorities for improvement of the framework. 